\documentclass[12pt]{article}

\usepackage[a4paper, total={6in, 8in}]{geometry}
\usepackage{textcomp}

\begin{document}
\setlength{\parindent}{0pt}

\def \t {\textrightarrow}
\def \v {\vspace{1em}}
\def \bi {\begin{itemize}}
\def \ei {\end{itemize}}
\def \s[#1] {\section*{#1}}
\def \ss[#1] {\subsection*{#1}}
\def \sss[#1] {\subsubsection*{#1}}

\s[Polisaccaridi]
L'amido si trova nei polisaccaridi \\
Il glicogeno è l'equivalente per gli organismi eterotrofi \\
La differenza tra i lipidi e il glicogeno \t il glicogeno non è legato all'adipe, ma si accumula nel fegato che è l'oragano + reattivo \t è il centro di purificazione delle attivita metaboliche \\
Infatti è quello + soggetto alle malattie \\
Anche nei muscoli si accumula \t hanno bisogno di energia da impiegare \\
L'energia viene presa dalle cellule \t quindi il glicogeno deve essere vicnio alle cellule che ne hanno bisogno \t è quindi una riserva

\v

La chitina \t nell'esosceletro degli insetti/crostacei \t diversa dalla cheratina, che è nei capelli e nelle unghie \\
La cellulosa è una molecola lineare \t ma da sola non basta \t se ci fosse solo lei, sarebbe elastica \t ha bisogno di un altro polisaccaride, che fa da reticolo e permette una struttura solida

\v

I monosaccaridi sono semplici \\
Noi non ci nutriamo di glucosio in realta \t non lo trovo infatti come monosaccaride, ma lo trovo associato ad altri saccaridi \\
I disaccaridi sono i + importanti \t hanno gruppi OH che a secondo della posizione (levo/destro) definiscono la funzione?? \\
Il legame O-glicosidico \t OH reagisce con l'altro ossidrile, c'è una condensazione \\
Ma cos'è un anomero? \t nei monosaccaridi c'è una forma aperta o una forma chiusa \t nel passaggio dalla forma aperta alla forma chiusa si determinano le caratteristiche del glucosio \\
C'è una coesistenza tra le due forme \t e l'ossidrile anomerico determina questo \t nella reazione di condensazione (ovvero la reazione che produce dell'acqua) reagisce l'OH e l'O con un ponte di ossgeno \\
Il legame si forma tra il gruppo caratterisitco e l'HO \t cosi avviene la ciclizzazione e l'OH è anomerico \\
Con $\alpha$ si indica il legame (carbonio $\alpha$ è il carbone vicino al legame aldedico/chetoso)

\v

Il lattosio è il disaccaride nel latto \t ha un glucosio e un galattosio \t il galattosio per essere digerito deve essere modificato in glucosio, mentre il glucosio viene subito metabolizato \\
Bisogna quindi prima idrolizzare \t enzima dell'elattomateri ?? permette di digerire il latte

\v

Omopolisaccaride = è un polisaccaride costituito dalla ripetizione di un singolo monosaccaride, e la catena può essere lineare o ramificata \\
Mentre eteropolisaccaride ha diversi monosaccaridi \\
L'amido è un omopolisaccaride \\
L'amido è costituito dall'amidosio e amilopectina \t non lo digeriamo direttamente perche non abbiamo l'enzima che accellera la separazione \\
Gli enzimi sono catalizzatori e sono sempre proteine qua \t il catalizzatore velocizza la reazione, e la rende efficace \t la reazione magari avviene, ma molto lentamente (come nel caso della digestione dell'amido) \\
Ma l'enzima non interviene nella reazione in se \t è un substrato che permette il verificarsi della reazione, ma non è presente nell'equazione chimica \t non interagisce chimicamente \\
Rendono + efficace l'urto \t e permettono anche di far avvenire la reazione a temperature più basse (abbassando la soglia energetica)

\s[Bo]
Isomeria ottica \t la molecola viene deviata in maniera diversa, perchè c'è un centro nella molecola \t un carbonio \\
Chiralita \t simmetria ma non sovrapponibilita \t puo accadere solo quando ho un carbonio legato a 4 costituenti diversi \t se solo due, allora si ha simmetria \\
La proiezione di Fischer \t quando si parla della enantiomeria si usano per indicare i 4 sostituenti diversi \\
Carbonio C1 è il carbonio + ossidato (ovvero ha carica parziale + positiva) \t quindi quello con il gruppo carbonilico e si mette in alto

\v

Reazione di ciclizzazione \t prevale la parte ciclica in soluzione acquosa, ed è la + stabile \\
Per metabolizzare il glucosio ho bisogno però la forma aperta \t e questa è la caratteristica importante del glucosio, e perche gli zuccheri sono usate per l'energia \\
Il D-glucosio è quello che metabolizziamo \t il L-glucosio invece no \t i nostri enzimi si attaccano al D

\ss[Formule di proiezione di Haworth]
Maggiore spessor eai legami che sono frontali \t si vuole far vedere la molecola ciclica inclinata \\
C1 e C5 etc. \t sono i legami che formano la catena, dei carboni ???

\v

Il carbonio anomerico è il carbonio che prima era aldedico mentre ora con la ciclizzazione si chiude e diventa stereogenico \t genera una stereoisomeria \\
I due isomeri si chiamano anomeri: alpha-d-glucosio e beta-d-glucosio

\v

I legami tra i monomeri sono da sapere \t nei casi dei monosaccaridi studiamo carboni aldosi e chetosi \\
Quando abbiamo polimeri, parliamo di oligosaccaridi (hanno pochi monomeri, sotto 10 circa) \\
Lattosio è un disaccaride \t composto da glucosio e galattosio \t vengono separati in due, e ci deve essere una catena che mi permette di tramutare ?? \\
Uomo si è evoluto per digerire il latte \t in realta grazie alla simbiosi con una flora batterica \\
Un gruppo ossidrilico reagisce con un altro gruppo ossidrilico \t si forma acqua (reazione di condensazione) e si forma il legame alpha1,4-glicosidico \t disaccaridi sono acetali, quindi reazione si chiama acetalizzazione (che è una reazione di condensazione)

\v

I polisaccaridi fanno da riserva energetica, sono strutturali (chitina e cellulosa) \\
Amido serve come riserva energetica, cellulosa è vegetale per struttura mentre chitina struttura ma animale \\
L'amidosio è elicoidale ma è una catena lineare fondamentalmente, mentre la amilopectina è ramificata \t e l'insieme dei due fornisce la struttura dell'amido \\
Il glicogeno è riserva energetica ma è animale \t nel fegato e nei muscoli, deve essere mobilitato velocemente \\
Il glucosio viene sintetizzato quando è in piu \t quando siamo in devito di energia, lo anabolizziamo (lo degradiamo) \\
Insulina e glucagone \t sono importanti per il diabete: senza insulina non riesco a sotrattre glucosio dall'organismo \\
Il pancreas produce e rilascia gli ormoni

\v

La cellulosa è un omopolisaccaride \t ma ha bisogno di un altro polisaccaride \\
Ha una forma lineare \t perche deve servire una funzione strutturale \\
La lignina ha una struttura a reticolo \t le fibre di cellulosa si incastrano nella lignina, che ha una forma a reticolo \t la cellulosa in se è flessibile, ma serve rigidita \\
Alcuni alberi sono + flessibili \t quelli che crescono in ambienti climatici forti, che devono proteggersi dal vento e da altri agenti atmosferici \\
La lignina permette alla cellulosa di organizzarsi in fasci e questo permette di avere rigidita \t diverse percentuali dei due determinano + o - flessibilita

\s[Lipidi]
Il glicogeno viene convertito in adipe \t che fa da riserva di energia, ma anche da isolante termico \\
La cellulosa è un omopolisaccaride \\
Vitamine non vengono sintetizzate dagli umani \\
I lipidi si dividono in saponificabili e non-saponificabili \t se tratto lipidi con una base forte, diventano sapone \\
I lipidi sono solubili in solventi apolari (etere, cloroformio etc.) e quindi non in acqua, che è polare \\
Saponific. e non sono molto diversi \\
Saturo o insaturo per il numero di doppi legami \\
I semplici \t sono esteri di acidi grassi con alcoli \\
Il fatto che sia grasso dipende dalla presenza del glicerolo \\
Gli eteri hanno un ponte ossigeno che lega i due carboni, e nasce dalla reazione di due gruppi ossidrilici \t ma se questa reazione si realizza con un cruppo carbossilico, porta alla creazione di un estere (ovvero -COO-) ?? \\
Mentre i lipidi complessi sono esteri di acidi grassi con alcoli che contengono gruppi atomici \t per esempio i fosfolipidi, ovvero quelli che costituiscono le membrane \\
Ci sono anche gli amminolipidi (gruppo $NH_2$, ovvero il gruppo amminico e ha carattere basico) \\
Infine, i "precursori e derivati lipidici" non sono nessuno dei due e si convertono in altre molecole \t es. sali biliari etc. \\
Acidita degli alcoli è bassa \\
Sono acidi grassi gli acidi carbossilici che hanno 4 o + carboni \t quelli con catena + lunga non sono solubili, perche la catena carboniosa + è lunga meno interagisce con l'acqua \t anche con gruppi elettronegativi???

\ss[Denominazione degli acidi grassi]
Il doppio legame ha una sua reattivita (è un addensamento di elettroni, quindi ha natura nucleofila, ed è forte e corto) \t e nel doppio legmane non è flessibile la catena, assume un angolazione precisa \\
È importante sapere la distanza tra il doppio legame e COOH \t se sono vicini si influenzano, se lontani no \\
Estremita metilica è quela senza COOH \t e chiamo omega il carbonio dell'estermita metilica (con $CH_3$, omega sta per ultimo perche è la ultima lettera dell alf.) \\
Non li possiamo produrre \t si chimano quindi grassi essenziali, non li possiamo produrre ma li dobbiamo assumere con la dieta 

\ss[Trigliceridi]
Sono dannosi ma sono anche necessari \t fanno da riserva energetica e negli animali formano il tessuto adiposo \t è il grasso che accumulo \\
Si chiamano invece oli nelle piante \t l'olio lo estraiamo e sono riserve energetiche \t all'oliva l'olio serve perche fa da riserva per il germoglio \\
Il trigliceride bisogna prima spezzarlo e ricavare i carboidrati \t per ricavare l'energia \\
1,2,3-propantriolo è un trigliceride \\
Il trigliceride si forma con il legame del glicerolo con tre acidi grassi \\
Una esterificazione \t reazione fra on OH e un COH porta a un composto che contiene O-C=O, ovvero un estere \\
Tre acidi grassi + glicerolo mi da trigliceride + acqua, perche estereficazione produce acqua \\
Si hanno catene lineari, emntre al centro si ha un doppio legame = geometria irregolare della molecola, quindi ha un certo ingombro e comportamenti diversi \\
I grassi animali sono ricchi di grassi saturi (senza doppi legami, e a temperatura ambiente sono solidi \t ma a temperatura elevata si sciolgono) \\
Invece l'olio a temperatura ambiente è liquido, è grasso insaturo

\ss[Saponificazione dei trigliceridi]
NaOH e KOH sono basi forti \t se si immerge in questi grassi caldi si formano i saponi \t si elimina il glicerolo e si ottiene sale \\
La molecola di sapone è una molecola anfipatica \t ovvero sia polare sia apolare \t infatti i saponi reagiscono con'acqua in cui si sciolgono, ma reagiscono con i grassi e infatti puliscono \\

\ss[Idrogenazione dei oli vegetali]
Gli oli sono insaturi \t e l'insaturazione porta a maggiore solubilita ??

\s[Lipidi]
Un lipide è caratterizzato da: (principale caratteristica) è insolubile in acqua \t quindi mancano legami che permettono legami a idrogeno con acqua \t quindi sono catene lunghe di carbonio \t sono composti apolari \\
Quelli grassi hanno COOH \\
Glicerolo si puo esterificare con tre diversi acidi grassi \\
Un acido grasso è caratterizzato dalla saturazione (grassi saturi non hanno doppi legami) \t e perchè è importante la presenza o meno dei doppi legami? \t il doppio legame è una zona nucleofila, e ha un impatto sulla forma \t un doppio legame irrigidisce \\
Questo provoca una precisa forma, che non è quella di una catena lineare che si può adattare \\
Il doppio legame cambia la forma perchè modifica una particolare configurazione: determina una geometria planare o una angolare, perchè cambia l'ibridazione (con il doppio legame è sp2)

\v

L'ibridazione è il fenomeno in cui avviene un rimodellamento della struttura esterna di un atomo, per cui ottegono degli orbitali intermedi tra gli atomi \\
Carbono ha numero atomico 6 \t quindi le scrittura degli orbitali è ?? \\
Mi aspetto nei diversi carboni con i diversi idrogeni degli orbitali diversi \t ma scopro che sono tutti uguali, quindi devono essere ibridati \\
Devo avere 75 per cento di p e 25 per cento di s \t sovrapposizione che permette di ottenere 4 orbitali che per 3/4 hanno natura p, 1/4 s \\
Uno di questo non è ibridato, un p rimane solo \t ho un orbitale p e due sp2

\v

Grassi animali sono saturi e solidi a temperatura ambiente \t perche? \t liquido o solido è determinato dalla forza attrativa delle molecole \\
Le forze di Van Der Waals \t prevedono che ci sono dei momenti in cui l'elettrone è spostato rispetto a dove ci si aspetta che sia, e questo cambia la configurazione \t sorge quindi una forza elettrica, e l'elettrone della molecola vicina si sposta etc. \\
Forze di Van Der Waals tengono insieme molecole che apparentemente sono neture \\
Quelle insature invece hanno gruppi negativi, che fanno pero anche da forze repulsiva \t insaturi provocano minori legami, e questo rende liquido a temperatura ambiente

\v

\ss[Saponificazione]
Reazione che avviene tra acidi grassi e una soluzione di base forte \t un ancido grasso ha gruppo COH, e con la base forte si forma un sale \\
I lipidi si dividono in saponificabili e non-saponificabili \\
Trigliceride con NaOH \t si ottiene uno ione ossalato che ?? \\
È una molecola anfipatica \t sia polare sia apolare \\
La margarina è solida ma è vegetale \t perché? \t è l'effetto dell'idrogenazione, ovvero l'aggiunta idrogeni a un doppio legame \t con l'idrogenazione posso ottenere la margarina, che è un olio che viene idrogenato \t cosi i doppi legami scompaiono e diventa solido \\
Si usa così perchè è meno ossidato rispetto ai doppi legami (quindi stabilizzato) e sostituisce il grasso animale del burro

\s[Lipidi]
Alifatico = catena di carboni legati da legami singoli \\
Triglicerdi sono composti da tre acidi grassi e un glicerolo \t e le catene di acidi grassi possono essere saturi o insaturi \\
Servono come riserva energetica

\ss[Fosfogliceride]
Un fosfogliceride \t glicerolo + 2 acidi grassi + fruppo fosfato \t il gruppo fosfato ha sempre una funzione energetica \\
ATP \t prodotta dai mitocondri ed è il risultato della respirazione cellulare \\
Il glicerolo esterifica \t ovvero forma un estere (gruppo funzionale O=C-O) \\
Testa idrofila \\
I fosfogliceridi formano le membrane cellulari \t è una molecola anfipatica, con una parte idrosolubile e una parte liposolubile ? \t ha una parte idrofoba e una idrofila

\ss[Terpeni]
Butadiene \t molecola con 4 carboni con due doppi legami \\
Terpeni sono isopropeni (2-metil-1,3 butadiene) \\
Ha una spiccata tendenza alla polimerizzazione \\
Non è saponificabile \t per saponificazione si ha bisogno di una base forte, e lui non ha una natura acida \t non puo reagire con una base forte \\

\ss[Colesterolo]
Ha 27 atomi di carbonio \t ed è una catena alchilica (ovvero senza doppi legami, alchili sono i gruppi funzionali che derivano dagli alcani) \\
Il colesterolo viene prodotto dal fegato \t e ne produce la quantita necessaria \\
Il colesterolo pero si introduce anche con la dieta \t e se ne assumo in eccesso, si creano problemi \\
Il colesterolo sballato puo esserci anche per disfunzioni epatiche \\
È importante perche se si mettono tra le code fosfolipidche cambiano la permeabilita della membrana \t influenzano quindi la fluidita delle membrane cellulari \\
Ne regolano questo aspetto \t membrane hanno diversi compiti: in alcune cellule devono essere + forti, in altre devono far passare molecole \t colesterolo regola la fluidita \\
Regola anche la temperatura corporea \t a temperatura ambiente (è stabile e ha cuore ciclico, è solido quindi) garantisce compattezza \\
A basse temperatura impedisce alla membrana di cristallizzarsi \t a basse temp. c'è il rischio che si irrigidisca, ma colesterolo previene questo \\
Colesterolo è importante anche per gli steroidi \t la vitaima D (composto che non puo essere sintentizzato) si produce con il sole \\
La bile \t nel fegato permette depurazione

\v

Se voglio digerire un grasso, lo emulsiono e cosi aumento superifcie di contatto \t cosi enizimi agiscono + efficientemente \t i sali biliari fanno questo

\s[Fosfogliceridi]
Hanno due acidi grassi, e poi un fosfato legato un altra catena (di solito è un amminoalcol) \\
Estereficazione si crea il legame COO \t COH + OH \t condensazione, si elimina l'aqua \\
Acido fosfatifico \t importante nei legami forti, si vede anche nell ATP \\
Testa idrofila \t carica negativa ? che porta a creazione di legami a idrogeno \\
Teste polari sono a contatto con il liquido \t mentre le code idrofobiche all opposto

\s[Terpeni]
2-metil-1,3-butandiene (diene perchè ha due doppi legami) \t questo è l'isoprene \\
È una molecola instabile \t tende alla polimerizzazione

\s[Colesterolo]
I saponi sono dei sali a catena lunga \\
I terpeni non sono saponificabili \t perchè non c'è la parte che puo formare il sale: isoprene non ha acidi grassi, non ha la parte che puo formare il sale \\
Anche il colesterolo \t 3 esagoni e un pentagono \t è uno steoride = gruppo di lipidi saponoficabili con idrucarburo policiclico saturo (saturo ovvero presenta doppio legame) \\
Alchile = catena dove i carboni sono legati con legami semplice \\
Colesterolo lo assumiamo ma lo produciamo anche \t e la produzione si regola con gli ormoni \t ma se lo introduco con alimentazione, non lo controllo \\
Colesterolo determina fluidita della membrana cellulare \t e ogni cellula ha sua funzionalita per cui serve una diversa fluiditia \t il colesterolo regola questo \\
Reazioni anabolizzanti = reazioni che costruiscono \t quando digerisco invece è reazione catabolizzatrice (distruggo) \\
Quando ho bisogno di colesterolo, chiedo al fegato di anabolizzare \\
Inoltre colesterolo impedisce che la membrana cristallizzi

\v

Sali biliari \t bile serve ad emulsionare i grassi = fenomeno per cui i lipidi vengono separati in piccole porzioni che vengono isolate tra loro da queste molecole, e questo rende + trasportabili e digeribili

\v

Ormoni sono idrofili \t mentre ormoni tiroidei e quelli delle gonadi sono lipofili \\
Cortisone ha una funzione antifiammatoria

\s[Recuperare 24 febbraio]
\s[Proteine]
Cos'è una proteina? \t è un polimero, il cui monomero è l'amminoacido che ha gruppo carbossilico (COOH) e ammidico \\
Si ha stereocentro di simmetria chirale in corrispondenza del carbonio $\alpha$ \\
Ci sono legami peptidici, che coinvolgono un gruppo carbossilico e un gruppo ammidico di due diversi amminoacidi \\
Le proteine si differenziano per la combinazione di amminoacidi, che sono 20 \\
Gli enzimi sono proteine, che fanno da catalizzatori \\
Sapere le varie funzioni: catalitica, di trasporto, strutturale, difesa, segnalazione, etc. \\
Proteina semplice è fatta solo di amminoacidi \\
Proteina coniugata = ha gruppi prostetici, come carboidrati o altre molecole organiche, o ioni metallici \\
Proteina fibrosa ha funzione strutturale: ha struttura lineare assembabile, mentre una proteina globulare no: non ha compatezza \\
Proteine fibrose sono insolubili, ma possono esserlo a condizioni particolari (alte temperature etc.) \t es. cheratina \\
Mioglobina è globulare, porta ossigeno nei muscoli e quindi deve essere solubile

\v

Polarita determina la solubilita in acqua \t pero cambia sempre in base alla lunghezza della catena idrofobica, che puo contrastare il gruppo elettronegativo \\
Zwitterione \t gruppo NH2 e COOH \t l'amminoacidi sono presenti nella stessa molecola, quindi ha comportamento acido/base, che dipenderà dal ph dell'ambiente (se ho ambiente acido, molecola si comportera come base ??) \\
In questo caso i valori del ph sono abbastanza neutri, quindi ho uno ione dipolare \\
\textbf{Definizione di acido: è un composto che in soluzione acquosa rilascia ioni H+} \t ambiente acido è ambiente con tanti ioni H+ \\
Quindi ioni H+ si legeranno + probabilmente alla parte basica con lo zwitterione

\v

Punto isoelettrico \t che è spesso intorno alla neutralita (ph 7) \\
La prolina è apolare, perche i due gruppi + e - si neutralizzano \t se la catena laterale è neutra, allora è apolare (alifatico non è aromatico) \\
Con catena aromatica invece è apolare perchè è molto stabile \\
Gruppi R carichi negativamente allora sono acidi, R positivi basici

\v

Perche parliamo di amminoacidi essenziali? \t con questi si indicano gli amminoacidi che servono al corpo umano ma non li puo sintetizzare \t e sono 8 \\
Con la dieta dobbiamo introdurre le vitamine e gli amminoacidi essenziali \\
La cisteina altera la proteina \t e la reazione però è reversibile

\s[Proteine]
Nel RER (reticolo endoplasmatico ruvido) \\
Gruppi prostetici sono cose che si aggiungono \\
Emoglobina ha il gruppo eme \t il gruppo prostetico è il ferro per ossigeno \t l'anemia falciforme è resistente alla malaria \\
Le proteine si classificano anche per la loro forma \t perche ne determina la funzione \\
\bi
    \item proteine lineari = funzione strutturale, per esempio $\alpha$-cheratina, che compone unghie e capelli (o il corno del rinoceronte) 
    \item proteine globulari = sono quelle trasportate, perche si attaccano \t + imp. è la mioglobina, che lega l'ossigeno ai muscoli
\ei

\v

Gli amminoacidi sono migliaia, ma quelli che compongono le nostre proteine sono 20 e sono $\alpha$-amminoacidi (per il carbonio $\alpha$, che lega COH a NH2) \\
Gruppi fond.: gruppo carbossilico (COOH) e gruppo amminico (NH2) \\
Da un amminoacido all altro cambia la catena laterale \t ogni amminoacido ha un centro chirale \\
L'amminoacido si presenta in ambiente neutro come dipolare (con entrambe le cariche) \t cambia in ambiente basico o acido \\
Punto isoelettrico \t un ponto in cui l'amm risulta neutro (elettricamente)

\v

AUG è mRNA, mentre TAC è DNA \\
Le proteine devono essere reattive, ma anche esse stabili \t quelle denaturate sono sono più funzionali
\end{document}
